\chapter{Introduction}

This blueprint follows the paper \emph{Combinatorial curve neighborhoods for
the affine flag manifold of type $A_1^{(1)}$}.  The formalization works in the
infinite dihedral group $D_\infty$ and proves the paper's main combinatorial
formula
\[
  \Gamma_d(u)=\{u w \mid w\in \max A_d(u)\}.
\]
The Lean development keeps the paper-level objects visible: chains in the
moment graph, the degree map $\varphi$, the auxiliary set $A_d(u)$, the maximal
root $\alpha(d)$, and the curve neighborhood $\Gamma_d(u)$.

\chapter{Section 2A: The Infinite Dihedral Group}

\begin{theorem}[Generators]\label{gen_by_sr0_sr1}
  \lean{gen_by_sr0_sr1}
  \leanfile{Dihedral/Basic.lean}
  \leanok
  The group $D_\infty$ is generated by the two simple reflections $sr(0)$ and
  $sr(1)$.
\end{theorem}

\begin{definition}[Length]\label{length_def}
  \leanfile{Dihedral/Basic.lean}
  \leanok
  Let $\ell : D_\infty \to \mathbb N$ be the Coxeter length.
\end{definition}

\begin{theorem}[Reduced-word length]\label{length_eq}
  \lean{length_eq}
  \leanfile{Dihedral/ReducedWord.lean}
  \uses{length_def}
  \leanok
  For every $g\in D_\infty$,
  \[
    \ell(g)=\operatorname{length}(\mathrm{reducedWord}(g)).
  \]
\end{theorem}

\begin{theorem}[Reflection length]\label{length_sr_abs}
  \lean{length_sr_abs}
  \leanfile{Dihedral/ReducedWord.lean}
  \uses{length_eq}
  \leanok
  For every $k\in\mathbb Z$,
  \[
    \ell(sr(k))=|2k-1|.
  \]
\end{theorem}

\begin{lemma}[Length dichotomy]\label{length_mul_eq_add_or_sub}
  \lean{length_mul_eq_add_or_sub}
  \leanfile{Dihedral/Degree.lean}
  \uses{length_def}
  \leanok
  If $\ell(u)\le \ell(v)$, then
  \[
    \ell(uv)=\ell(u)+\ell(v)
    \quad\hbox{or}\quad
    \ell(uv)=\ell(v)-\ell(u).
  \]
\end{lemma}

\begin{definition}[Roots and root reflections]\label{root_def}
  \lean{Root, s_α, length_root_reflection, φ_s_alpha_eq}
  \leanfile{Dihedral/Degree.lean}
  \leanok
  A positive root is a pair $(a,b)\in\mathbb N^2$ with $a=b+1$ or $b=a+1$.
  Its root reflection is the alternating word $s_\alpha$ of length $a+b$, and
  Lean proves $\ell(s_\alpha)=a+b$ and $\varphi(s_\alpha)=\alpha$.
\end{definition}

\begin{definition}[Degree of an element]\label{degree_def}
  \lean{Degree, getDegree}
  \leanfile{Dihedral/Degree.lean}
  \leanok
  A degree is a pair of natural numbers with product order.  The map
  $\varphi : D_\infty\to\mathbb N^2$ records the numbers of $s_0$ and $s_1$ in
  the reduced word.
\end{definition}

\chapter{Section 2B: Moment Graph, Chains, and Order}

\begin{definition}[Edges and chains]\label{chain_def}
  \lean{IsEdge, HasChain, HasIncreasingChain}
  \leanfile{Dihedral/Chain.lean}
  \uses{root_def, degree_def}
  \leanok
  An edge $u\xrightarrow{\alpha}v$ means $v=u s_\alpha$.  A chain is a finite
  succession of such edges; its degree is the sum of the root degrees on the
  edges.
\end{definition}

\begin{theorem}[Bruhat order by length]\label{lt_iff_length_lt}
  \lean{lt_iff_length_lt}
  \leanfile{Dihedral/Chain.lean}
  \uses{chain_def, length_def}
  \leanok
  For vertices $u,v$,
  \[
    u<v \iff \ell(u)<\ell(v).
  \]
\end{theorem}

\begin{lemma}[Paper Lemma 2.5]\label{lemma_2_5}
  \lean{lemma_2_5_a, lemma_2_5_b, chain_degree_parity, inv_mul_degree_le_of_chain}
  \leanfile{Dihedral/Chain.lean}
  \uses{chain_def, degree_def}
  \leanok
  If there is a chain from $u$ to $v$ of total degree $d$, then
  \[
    d=\varphi(u^{-1}v)+2(r,s)
  \]
  for some $r,s\in\mathbb N$.  In particular $\varphi(u^{-1}v)\le d$.
\end{lemma}

\begin{definition}[Reachable set and curve neighborhood]\label{curve_nbhd_def}
  \lean{ReachableSet, CurveNeighborhood, IsMaximalIn}
  \leanfile{Dihedral/Ad.lean}
  \uses{chain_def}
  \leanok
  The reachable set is
  \[
    \mathrm{ReachableSet}(u,d)=\{v\mid \exists d'\le d,\;\texttt{HasChain }u\ v\ d'\}.
  \]
  The curve neighborhood $\Gamma_d(u)$ is the set of maximal elements of this
  reachable set.
\end{definition}

\chapter{Section 3: Calculation of Curve Neighborhoods}

\begin{definition}[The set $A_d(u)$]\label{ad_def}
  \lean{Ad}
  \leanfile{Dihedral/Ad.lean}
  \uses{degree_def, length_def}
  \leanok
  For $u\in D_\infty$ and degree bound $d$,
  \[
    A_d(u)=\{v\mid \ell(uv)=\ell(u)+\ell(v),\ \varphi(v)\le d\}.
  \]
\end{definition}

\begin{definition}[The maximal root $\alpha(d)$]\label{max_root_le_def}
  \lean{maxRootLE, root_from_degree, rootReflectionLE, s_alpha_d}
  \leanfile{Dihedral/Identity.lean}
  \uses{root_def, degree_def}
  \leanok
  For non-diagonal $d$, $\alpha(d)$ is the maximal root below $d$:
  \[
    \alpha(d)=
    \begin{cases}
      (d.a,d.a+1), & d.a<d.b,\\
      (d.b+1,d.b), & d.a>d.b.
    \end{cases}
  \]
  Lean keeps the old names \texttt{root\_from\_degree} and \texttt{s\_alpha\_d},
  and exposes the semantic aliases \texttt{maxRootLE} and
  \texttt{rootReflectionLE}.
\end{definition}

\begin{lemma}[Paper Lemma 3.1]\label{lemma_3_1}
  \lean{exists_unique_max_ad_ne_one, exists_unique_max_ad_one_neq, max_ad_one_diag_iff_r_or_neg_r}
  \leanfile{Dihedral/Ad.lean}
  \uses{ad_def}
  \leanok
  The set $A_d(u)$ has a unique maximal element unless $u=1$ and $d=(a,a)$.
  In the exceptional diagonal identity case, the maximal elements are
  $r(a)$ and $r(-a)$.
\end{lemma}

\begin{theorem}[Formula (3): identity base point]\label{formula_3_identity}
  \lean{curve_nbhd_one, curve_nbhd_one_diag, curve_nbhd_one_offdiag_rootReflectionLE}
  \leanfile{Dihedral/Identity.lean}
  \uses{lemma_3_1, curve_nbhd_def, max_root_le_def}
  \leanok
  At the identity,
  \[
    \Gamma_d(1)=
    \begin{cases}
      \{(s_0s_1)^a,(s_1s_0)^a\}, & d=(a,a),\\
      \{s_{\alpha(d)}\}, & d.a\ne d.b.
    \end{cases}
  \]
\end{theorem}

\begin{theorem}[Formula (4): one endpoint case]\label{formula_4_ends_s0}
  \lean{maxAdCandidateEndsS0, max_ad_candidate_ends_s0, max_ad_ends_s0}
  \leanfile{Dihedral/Ends.lean}
  \uses{ad_def, max_root_le_def}
  \leanok
  If the reduced word for $u$ ends in $s_0$, Lean defines the paper's candidate
  maximal element by \texttt{maxAdCandidateEndsS0 d}.  The theorem
  \texttt{max\_ad\_candidate\_ends\_s0} proves that this candidate is maximal
  in $A_d(u)$.  The case ending in $s_1$ is symmetric and is not duplicated as a
  separate public formula.
\end{theorem}

\begin{lemma}[Paper Lemma 3.4]\label{lemma_3_4}
  \lean{lemma_3_4, len_mul_ad_le_curve_nbhd_max, inv_mul_len_le_curve_nbhd_one_max}
  \leanfile{Dihedral/CurveNeighborhood.lean}
  \uses{curve_nbhd_def, ad_def, formula_3_identity}
  \leanok
  If $z\in A_d(1)$ and $v\in\Gamma_d(u)$, then $\ell(uz)\le\ell(v)$ and
  $\varphi(u^{-1}v)\le d$.  If moreover $z\in\Gamma_d(1)$, then
  $\ell(u^{-1}v)\le\ell(z)$.
\end{lemma}

\begin{lemma}[Paper Lemma 3.5]\label{lemma_3_5}
  \lean{lemma_3_5, inv_mul_mem_ad_curve_nbhd}
  \leanfile{Dihedral/CurveNeighborhood.lean}
  \uses{lemma_3_4, ad_def}
  \leanok
  If $v\in\Gamma_d(u)$, then
  \[
    u^{-1}v\in A_d(u).
  \]
\end{lemma}

\begin{theorem}[Paper Theorem 3.6]\label{curve_nbhd_eq_mul_max_ad}
  \lean{theorem_3_6, curve_nbhd_eq_mul_max_ad, main_theorem}
  \leanfile{Dihedral/CurveNeighborhood.lean}
  \uses{lemma_3_5, ad_def, curve_nbhd_def}
  \leanok
  For every $u\in D_\infty$ and every degree bound $d$,
  \[
    \Gamma_d(u)=\{v\mid \exists w,\ w\in\max A_d(u)\ \hbox{and}\ v=u w\}.
  \]
\end{theorem}

\chapter{Computable Appendix}

\begin{definition}[Finite enumeration]\label{enumerateD_def}
  \lean{cLength, enumerateD, Ad_finset}
  \leanfile{Dihedral/Computable.lean}
  \uses{length_eq, ad_def}
  \leanok
  The computable model enumerates all elements up to the finite length bound
  $d.a+d.b+1$ and filters by the defining conditions of $A_d(u)$.
\end{definition}

\begin{theorem}[Finite model correctness]\label{computable_curve_nbhd_iff}
  \lean{mem_CurveNeighborhood_computable_iff, coe_CurveNeighborhood_computable, coe_Ad_finset}
  \leanfile{Dihedral/Computable.lean}
  \uses{enumerateD_def, curve_nbhd_eq_mul_max_ad}
  \leanok
  For every base point $u$, degree bound $d$, and vertex $v$,
  \[
    v\in\texttt{CurveNeighborhood\_computable }u\ d
    \iff
    v\in\Gamma_d(u).
  \]
\end{theorem}

\chapter{Local Support Lemmas}

\begin{lemma}[Integer and length API]\label{local_mathlib_api}
  \lean{natAbs_sub_one_of_pos, natAbs_sub_one_of_nonpos,
    natAbs_add_eq_natAbs_add_two_mul, natAbs_sub_eq_natAbs_add_two_mul,
    getDegree_r_neg_natCast}
  \leanfile{Dihedral/Mathlib.lean}
  \leanfile{Dihedral/Degree.lean}
  \leanok
  The local support files stage the integer absolute-value and rank-two length
  facts needed to keep the infinite-dihedral model stable under mathlib changes.
\end{lemma}